% !TeX program = xelatex
\documentclass[10pt,a4paper]{article}

% Korean/Unicode support: set Overleaf compiler to XeLaTeX.
\usepackage{kotex}
\usepackage[margin=24mm]{geometry}
\usepackage{amsmath,amssymb}
\usepackage{booktabs,longtable,array,calc,etoolbox,tabularx}
\usepackage{graphicx}
\usepackage{tikz}
\usetikzlibrary{arrows.meta,positioning,shapes.geometric}
\usepackage{xcolor}
\usepackage{hyperref}
\usepackage{bookmark}
\usepackage{xurl}
\usepackage{fancyvrb}
\usepackage{fvextra}
\usepackage{enumitem}
\usepackage{titlesec}
\usepackage{setspace}
\usepackage{caption}

\hypersetup{
  pdftitle={LENS - DACON 월간 해커톤 제출용},
  pdfauthor={최문석, 최예윤},
  colorlinks=true,
  linkcolor=black,
  urlcolor=blue,
  citecolor=black
}

\setstretch{1.12}
\setlength{\parindent}{0pt}
\setlength{\parskip}{6pt plus 2pt minus 1pt}
\setlength{\emergencystretch}{3em}
\setlength{\tabcolsep}{4pt}
\renewcommand{\arraystretch}{1.2}
\newcolumntype{Y}{>{\raggedright\arraybackslash}X}
\AtBeginEnvironment{longtable}{\footnotesize}

\titleformat{\section}{\Large\bfseries}{\thesection}{0.75em}{}
\titleformat{\subsection}{\large\bfseries}{\thesubsection}{0.75em}{}
\titleformat{\subsubsection}{\normalsize\bfseries}{\thesubsubsection}{0.75em}{}

\DefineVerbatimEnvironment{Highlighting}{Verbatim}{commandchars=\\\{\},breaklines=true,fontsize=\small}
\DefineVerbatimEnvironment{verbatim}{Verbatim}{breaklines=true,fontsize=\small}
\newenvironment{Shaded}{\vspace{0.3em}\begin{quote}}{\end{quote}\vspace{0.3em}}
\newcommand{\VerbBar}{|}
\newcommand{\VERB}{\Verb[commandchars=\\\{\}]}
\newcommand{\AlertTok}[1]{\textcolor[rgb]{1.00,0.00,0.00}{\textbf{#1}}}
\newcommand{\AnnotationTok}[1]{\textcolor[rgb]{0.38,0.63,0.69}{\textbf{\textit{#1}}}}
\newcommand{\AttributeTok}[1]{\textcolor[rgb]{0.49,0.56,0.16}{#1}}
\newcommand{\BaseNTok}[1]{\textcolor[rgb]{0.25,0.63,0.44}{#1}}
\newcommand{\BuiltInTok}[1]{\textcolor[rgb]{0.00,0.50,0.00}{#1}}
\newcommand{\CharTok}[1]{\textcolor[rgb]{0.25,0.44,0.63}{#1}}
\newcommand{\CommentTok}[1]{\textcolor[rgb]{0.38,0.63,0.69}{\textit{#1}}}
\newcommand{\CommentVarTok}[1]{\textcolor[rgb]{0.38,0.63,0.69}{\textbf{\textit{#1}}}}
\newcommand{\ConstantTok}[1]{\textcolor[rgb]{0.53,0.00,0.00}{#1}}
\newcommand{\ControlFlowTok}[1]{\textcolor[rgb]{0.00,0.44,0.13}{\textbf{#1}}}
\newcommand{\DataTypeTok}[1]{\textcolor[rgb]{0.56,0.13,0.00}{#1}}
\newcommand{\DecValTok}[1]{\textcolor[rgb]{0.25,0.63,0.44}{#1}}
\newcommand{\DocumentationTok}[1]{\textcolor[rgb]{0.73,0.13,0.13}{\textit{#1}}}
\newcommand{\ErrorTok}[1]{\textcolor[rgb]{1.00,0.00,0.00}{\textbf{#1}}}
\newcommand{\ExtensionTok}[1]{#1}
\newcommand{\FloatTok}[1]{\textcolor[rgb]{0.25,0.63,0.44}{#1}}
\newcommand{\FunctionTok}[1]{\textcolor[rgb]{0.02,0.16,0.49}{#1}}
\newcommand{\ImportTok}[1]{\textcolor[rgb]{0.00,0.50,0.00}{\textbf{#1}}}
\newcommand{\InformationTok}[1]{\textcolor[rgb]{0.38,0.63,0.69}{\textbf{\textit{#1}}}}
\newcommand{\KeywordTok}[1]{\textcolor[rgb]{0.00,0.44,0.13}{\textbf{#1}}}
\newcommand{\NormalTok}[1]{#1}
\newcommand{\OperatorTok}[1]{\textcolor[rgb]{0.40,0.40,0.40}{#1}}
\newcommand{\OtherTok}[1]{\textcolor[rgb]{0.00,0.44,0.13}{#1}}
\newcommand{\PreprocessorTok}[1]{\textcolor[rgb]{0.74,0.48,0.00}{#1}}
\newcommand{\RegionMarkerTok}[1]{#1}
\newcommand{\SpecialCharTok}[1]{\textcolor[rgb]{0.25,0.44,0.63}{#1}}
\newcommand{\SpecialStringTok}[1]{\textcolor[rgb]{0.73,0.40,0.53}{#1}}
\newcommand{\StringTok}[1]{\textcolor[rgb]{0.25,0.44,0.63}{#1}}
\newcommand{\VariableTok}[1]{\textcolor[rgb]{0.10,0.09,0.49}{#1}}
\newcommand{\VerbatimStringTok}[1]{\textcolor[rgb]{0.25,0.44,0.63}{#1}}
\newcommand{\WarningTok}[1]{\textcolor[rgb]{0.38,0.63,0.69}{\textbf{\textit{#1}}}}
\providecommand{\tightlist}{\setlength{\itemsep}{0pt}\setlength{\parskip}{0pt}}


\begin{document}
\begin{titlepage}
\centering
\vspace*{28mm}
{\Huge\bfseries LENS\par}
\vspace{6mm}
{\Large Skills.md 기반 AI 금융 포트폴리오 대시보드 서비스 개발 기획서\par}
\vspace{10mm}
{\large DACON 월간 해커톤 : 투자 데이터를 시각화하라- Skills 기반 대시보드 설계 제출용\par}
\vspace{24mm}
{\large
최문석\textsuperscript{1} \quad 최예윤\textsuperscript{2}\par}
\vspace{5mm}
{\normalsize
\textsuperscript{1}팀장, Data Scientist, 경북대학교 식품자원경제학과\par
\textsuperscript{2}팀원, Frontend Developer, 경북대학교 컴퓨터학부 글로벌소프트웨어 전공\par}
\vfill
{\large 2026.04.30.\par}
\end{titlepage}
\tableofcontents
\newpage
\section{Executive Summary}\label{executive-summary}

\textbf{공개용 한 줄 소개:}
``제각각인 투자 CSV를 규칙 기반 포트폴리오 대시보드로 바꾸는 AI 금융 분석 서비스.''

\textbf{기획서용 핵심 정의:}
LENS는 사용자의 투자 CSV/XLSX와 분석 목적을 함께 해석하여, Skills.md에 정의된 금융 분석 규칙에 따라 가능한 분석만 자동 실행하고, 재현 가능한 포트폴리오 대시보드와 AI 애널리스트 리포트를 생성하는 서비스다.

\textbf{대상 사용자:} 포트폴리오 성과 보고를 담당하는 금융 데이터 실무자(1차 타깃), 개인 퀀트 연구자 및 투자 분석 학습자(2차 타깃).

\textbf{공모전 적합성:} 이 공모전의 핵심 요구사항은 ``분석 규칙을 Skills.md에 문서화하고 그 규칙이 대시보드 생성에 실제로 작동함을 보여주는 것''이다. LENS는 Skills.md를 단순 설명 문서가 아닌 분석 파이프라인 전체의 판단 기준으로 설계하며, Skills 로그가 규칙-결과 연결을 기록한다.

\textbf{P0 데모 핵심 장면:} 비정형 한국어 CSV 파일을 드롭 → AI가 컬럼 자동 매핑 → 사용자 목적 확인 → 데이터 충분성 판정 → 누적 수익률·MDD·변동성·비중 대시보드 자동 생성 → AI 애널리스트 패널 표시 → Skills 로그로 규칙 추적.

\textbf{Skills.md의 역할:} Skills.md는 컬럼 매핑 추론 기준, 데이터 충분성 판정 조건, 지표 계산 공식, 시각화 차트 선택 기준, 인사이트 해석 템플릿, 금지 표현 검증 규칙, 재현성 계약을 모두 사전 정의한다. Skills 로그는 이 규칙들이 실제로 대시보드 생성에 적용되었음을 기록한다.

\begin{center}\rule{0.5\linewidth}{0.5pt}\end{center}

\section{공모전 요구사항 해석}\label{uxacf5uxbaa8uxc804-uxc694uxad6cuxc0acuxd56d-uxd574uxc11d}

\subsection{공모전의 본질적 요구}\label{uxacf5uxbaa8uxc804uxc758-uxbcf8uxc9c8uxc801-uxc694uxad6c}

공모전의 공식 명칭은 ``월간 해커톤: 투자 데이터를 시각화하라 --- Skills 기반 대시보드 설계''다. 공식 설명에 따르면 Skills.md는 투자 데이터 분석 기준, 투자 지표 계산 규칙, 시각화 선택 기준, 리포트 구성 흐름, 인사이트 생성 규칙을 포함하는 ``분석 규칙과 생성 지침을 담은 범용 MD 문서''를 의미한다.

본 기획서는 공모전의 핵심 평가 기준을 다음과 같이 해석한다. ``분석 규칙을 Skills.md에 문서화하고, 그 규칙이 실제 대시보드 생성에 적용되는 구조를 보여주는 것.'' 분석 규칙과 대시보드 출력 간의 연결 구조가 평가의 주된 기준이라고 판단한다.

\subsection{평가 항목 해석}\label{uxd3c9uxac00-uxd56duxbaa9-uxd574uxc11d}

\begin{longtable}[]{@{}
  >{\raggedright\arraybackslash}p{(\columnwidth - 4\tabcolsep) * \real{0.3600}}
  >{\raggedright\arraybackslash}p{(\columnwidth - 4\tabcolsep) * \real{0.2000}}
  >{\raggedright\arraybackslash}p{(\columnwidth - 4\tabcolsep) * \real{0.4400}}@{}}
\toprule\noalign{}
\begin{minipage}[b]{\linewidth}\raggedright
평가항목
\end{minipage} & \begin{minipage}[b]{\linewidth}\raggedright
배점
\end{minipage} & \begin{minipage}[b]{\linewidth}\raggedright
우리의 해석
\end{minipage} \\
\midrule\noalign{}
\endhead
\bottomrule\noalign{}
\endlastfoot
범용성 & 25 & 정의된 CSV/XLSX 데이터 유형(일별 수익률, 일별 가격, 보유 비중, 거래 내역) 내에서 컬럼 구조가 달라도 처리할 수 있는 유연성 \\
Skills.md 설계 & 25 & 규칙 ID 체계, 계산 공식, 임계값, 시각화 선택 기준, 인사이트 템플릿이 명시된 명세서 수준의 Skills.md \\
대시보드 자동 생성 & 25 & 업로드에서 대시보드 생성까지 수작업 없이 완주되는 흐름, Skills.md 규칙에 따른 자동 분석 \\
바이브코딩 활용 & 15 & Skills.md를 기준 문서로 삼아 AI 개발 도구가 화면·기능·분석 규칙 산출물을 일관되게 생성하고, 리뷰 에이전트가 P0 범위·금지 표현·재현성 조건을 검증하는 방식으로 바이브코딩을 활용한다 \\
실용성 및 창의성 & 10 & 데이터 충분성 검사 + 부분 분석 + Skills 로그 가시화의 실무 활용 가능성 \\
\end{longtable}

상위 3개 항목(범용성·Skills.md 설계·대시보드 자동 생성)의 합산이 75점이다. LENS의 설계는 이 세 항목의 연결고리인 ``Skills.md → 파이프라인 → 대시보드 → 로그''를 단일 구조로 구현한다.

\subsection{LENS의 해석}\label{lensuxc758-uxd574uxc11d}

본 기획서의 설계 목표는 금융 포트폴리오 분석 규칙을 Skills.md에 재사용 가능한 형태로 정의하고, 그 규칙이 대시보드 생성 과정 전체를 제어하는 구조를 기획 수준에서 명세하는 것이다.

\begin{center}\rule{0.5\linewidth}{0.5pt}\end{center}

\section{문제 정의}\label{uxbb38uxc81c-uxc815uxc758}

\subsection{문제 배경}\label{uxbb38uxc81c-uxbc30uxacbd}

금융 기업과 투자 실무 현장에서 포트폴리오 데이터는 Excel 또는 CSV 형식으로 수수하는 관행이 지속되고 있다. 거래내역 정산, 운용 성과 집계, 위험 지표 산출이 모두 이 파일 형식에서 출발하는 경우가 많다. 세 가지 변화가 이 문제를 가시화했다.

\textbf{첫째,} Julius AI, Power BI Copilot 등 AI 기반 데이터 분석 도구가 보급되었으나, 금융 도메인 분석 규칙은 사용자가 매 세션마다 프롬프트로 재정의해야 한다. Microsoft 공식 문서는 Copilot 생성 보고서가 비결정적(nondeterministic)이며 도메인 전문가를 대체할 수 없다고 명시한다.

\textbf{둘째,} Portfolio Visualizer, Sharesight, Koyfin 등 금융 특화 서비스는 성숙한 분석 기능을 제공하지만, 공개된 제품 설명을 검토하면 티커 기반 입력, 브로커 연동, 또는 자체 포트폴리오 등록 방식을 주된 데이터 진입 경로로 제시한다. 사용자의 비정형 CSV 컬럼명이 이들 시스템의 기대 스키마와 다를 경우 별도의 데이터 변환 작업이 필요하다 (기획 해석).

\textbf{셋째,} Skills.md 개념이 등장하면서, 분석 규칙을 사전에 문서화하고 그 문서에 기반해 대시보드를 자동 생성하는 구조가 기술적으로 가능해졌다. 이 구조를 공개된 제품 사용 흐름에서 전면에 제시하는 서비스는 확인되지 않는다 (기획 해석).

\subsection{핵심 문제 구조: User = Data + Goal}\label{uxd575uxc2ec-uxbb38uxc81c-uxad6cuxc870-user-data-goal}

사용자는 두 가지를 가지고 있다. 데이터(업로드한 파일)와 목적(분석하고 싶은 것)이다.

\begin{longtable}[]{@{}
  >{\raggedright\arraybackslash}p{(\columnwidth - 2\tabcolsep) * \real{0.5385}}
  >{\raggedright\arraybackslash}p{(\columnwidth - 2\tabcolsep) * \real{0.4615}}@{}}
\toprule\noalign{}
\begin{minipage}[b]{\linewidth}\raggedright
사용자가 가진 것
\end{minipage} & \begin{minipage}[b]{\linewidth}\raggedright
현재의 장벽
\end{minipage} \\
\midrule\noalign{}
\endhead
\bottomrule\noalign{}
\endlastfoot
포트폴리오 CSV (컬럼 구조 제각각) & 분석 서비스는 고정 스키마를 요구하는 경향이 있다 \\
분석 목적 (성과, 리스크, 비중 파악 등) & AI 도구는 매번 사용자가 규칙을 직접 정의해야 한다 \\
\end{longtable}

LENS는 이 두 가지를 연결하는 역할을 한다. 어느 한쪽이 불분명하면 분석을 강행하지 않고 명확화 단계를 먼저 수행한다.

\subsection{사용자 Pain Point}\label{uxc0acuxc6a9uxc790-pain-point}

\begin{longtable}[]{@{}
  >{\raggedright\arraybackslash}p{(\columnwidth - 6\tabcolsep) * \real{0.1071}}
  >{\raggedright\arraybackslash}p{(\columnwidth - 6\tabcolsep) * \real{0.3929}}
  >{\raggedright\arraybackslash}p{(\columnwidth - 6\tabcolsep) * \real{0.2143}}
  >{\raggedright\arraybackslash}p{(\columnwidth - 6\tabcolsep) * \real{0.2857}}@{}}
\toprule\noalign{}
\begin{minipage}[b]{\linewidth}\raggedright
\#
\end{minipage} & \begin{minipage}[b]{\linewidth}\raggedright
Pain Point
\end{minipage} & \begin{minipage}[b]{\linewidth}\raggedright
설명
\end{minipage} & \begin{minipage}[b]{\linewidth}\raggedright
페르소나
\end{minipage} \\
\midrule\noalign{}
\endhead
\bottomrule\noalign{}
\endlastfoot
P1 & 파일을 올려도 분석이 바로 시작되지 않는다 & 서비스 스키마와 사용자 CSV 컬럼명이 다르면 변환 수작업이 선행된다 & A, B \\
P2 & 같은 데이터를 넣어도 분석 결과가 달라진다 & 범용 AI 도구는 세션 간 재현성이 보장되지 않는다 & A \\
P3 & 데이터가 부족할 때 오류를 내거나 침묵한다 & 어떤 분석이 가능하고 무엇이 부족한지 안내받지 못한다 & A, B \\
P4 & 차트를 만들었지만 왜 그런 결과인지 설명하기 어렵다 & 계산 규칙과 판단 기준을 추적할 수 없다 & A, B \\
P5 & 분석 결과를 바로 보고서로 쓸 수 없다 & 차트·수치 별도 생성 후 수작업으로 문서에 붙여넣어야 한다 & A \\
P6 & 데이터 형식마다 분석 흐름을 다시 설계해야 한다 & 일별 수익률과 보유 비중 파일은 구조가 달라 매번 재구성이 필요하다 & A, B \\
\end{longtable}

\begin{center}\rule{0.5\linewidth}{0.5pt}\end{center}

\section{대상 사용자}\label{uxb300uxc0c1-uxc0acuxc6a9uxc790}

\subsection{페르소나 A: 투자 운용 실무자 (P0 1차 타깃)}\label{uxd398uxb974uxc18cuxb098-a-uxd22cuxc790-uxc6b4uxc6a9-uxc2e4uxbb34uxc790-p0-1uxcc28-uxd0c0uxae43}

\begin{longtable}[]{@{}
  >{\raggedright\arraybackslash}p{(\columnwidth - 2\tabcolsep) * \real{0.5000}}
  >{\raggedright\arraybackslash}p{(\columnwidth - 2\tabcolsep) * \real{0.5000}}@{}}
\toprule\noalign{}
\begin{minipage}[b]{\linewidth}\raggedright
항목
\end{minipage} & \begin{minipage}[b]{\linewidth}\raggedright
내용
\end{minipage} \\
\midrule\noalign{}
\endhead
\bottomrule\noalign{}
\endlastfoot
역할 & 포트폴리오 분석 담당자, 운용 성과 보고서 작성자, 금융 데이터 실무자 \\
보유 데이터 & 사내 시스템에서 내려받은 일별 수익률 CSV, 거래 상대방으로부터 수수한 Excel, 보유 종목 비중 시트 \\
분석 목적 & 월간·분기별 성과 보고: 누적 수익률, MDD, 연환산 변동성, 종목별 비중 \\
현재 도구의 한계 & Excel은 수식을 매번 재구성해야 하고 재현성이 낮다. 사내 보고 시스템은 정해진 양식에만 맞춰 사용해야 한다 \\
LENS 가치 & 비정형 CSV 그대로 업로드 → 자동 분석 → 규칙 기반 대시보드 + 애널리스트 리포트 \\
\end{longtable}

\subsection{페르소나 B: 개인 퀀트 연구자 (P0 2차 타깃)}\label{uxd398uxb974uxc18cuxb098-b-uxac1cuxc778-uxd000uxd2b8-uxc5f0uxad6cuxc790-p0-2uxcc28-uxd0c0uxae43}

\begin{longtable}[]{@{}
  >{\raggedright\arraybackslash}p{(\columnwidth - 2\tabcolsep) * \real{0.5000}}
  >{\raggedright\arraybackslash}p{(\columnwidth - 2\tabcolsep) * \real{0.5000}}@{}}
\toprule\noalign{}
\begin{minipage}[b]{\linewidth}\raggedright
항목
\end{minipage} & \begin{minipage}[b]{\linewidth}\raggedright
내용
\end{minipage} \\
\midrule\noalign{}
\endhead
\bottomrule\noalign{}
\endlastfoot
역할 & 퀀트 입문자, 금융 데이터 분석 학습자, 투자 분석 스터디 참여자 \\
보유 데이터 & 다운로드한 주가 CSV, ETF 수익률 데이터, 직접 구성한 비중 시트 \\
분석 목적 & 성과 분해, 드로다운 분석, 기초 리스크 지표, 비중 분포 확인 \\
현재 도구의 한계 & Python은 시각화·리포트 정리에 추가 작업이 필요하다. 범용 AI 도구는 금융 지표 정의를 매번 포함해야 한다 \\
LENS 가치 & 분석 규칙 재정의 없이 동일 품질의 대시보드 재현 \\
\end{longtable}

\begin{center}\rule{0.5\linewidth}{0.5pt}\end{center}

\section{제품 개념}\label{uxc81cuxd488-uxac1cuxb150}

\subsection{LENS란}\label{lensuxb780}

LENS는 정의된 포트폴리오 CSV/XLSX 데이터 유형을 처리하는 웹 기반 금융 대시보드 서비스다. Skills.md에 정의된 규칙이 분석 흐름 전체를 제어하며, AI는 규칙 범위 안에서 컬럼 매핑 추론과 자연어 해석 보조에 한정된다.

\textbf{본 MVP 서비스 범위:}
- Skills.md에 정의된 규칙 기반 포트폴리오 분석 서비스
- 데이터 충분성을 먼저 확인하고 가능한 분석만 실행하는 서비스
- 모든 분석 단계의 적용 규칙을 Skills 로그에 기록하는 서비스

\textbf{본 MVP 범위 외:}
- 범용 BI 도구 기능 (임의 데이터 유형에 대한 범용 처리는 제공하지 않는다)
- 자유형 AI 챗봇 기능 (비정형 자연어 질의응답은 제공하지 않는다)
- 투자 추천·미래 수익률 전망·인과 관계 단정 출력
- 자동매매·브로커 연동

\subsection{핵심 워크플로}\label{uxd575uxc2ec-uxc6ccuxd06cuxd50cuxb85c}

\begin{longtable}[]{@{}
  >{\raggedright\arraybackslash}p{(\columnwidth - 4\tabcolsep) * \real{0.2308}}
  >{\raggedright\arraybackslash}p{(\columnwidth - 4\tabcolsep) * \real{0.2308}}
  >{\raggedright\arraybackslash}p{(\columnwidth - 4\tabcolsep) * \real{0.5385}}@{}}
\toprule\noalign{}
\begin{minipage}[b]{\linewidth}\raggedright
단계
\end{minipage} & \begin{minipage}[b]{\linewidth}\raggedright
동작
\end{minipage} & \begin{minipage}[b]{\linewidth}\raggedright
Skills.md 관여
\end{minipage} \\
\midrule\noalign{}
\endhead
\bottomrule\noalign{}
\endlastfoot
1. 파일 업로드 & CSV/XLSX 수신, 형식 검증 & 지원 형식 목록 \\
2. 데이터 미리보기 & 상위 5행 + 컬럼 목록 표시 & --- \\
3. AI 컬럼 매핑 제안 & 컬럼명 → 의미 범주 추론 + 신뢰도 & RULE-CM-xxx \\
4. 사용자 매핑 확정 & 확인·수정·확정 & 확정 후 결정적 \\
5. 분석 목적 입력 & 자연어 또는 후보 카드 선택 & RULE-GC-xxx \\
6. 목적 모호성 명확화 & 모호한 입력 시 후보 제시 → 확인 & RULE-GC-xxx \\
7. 데이터 충분성 검사 & 항목별 가능/제한/불가 판정 & RULE-DS-xxx \\
8. 가능한 분석 자동 실행 & 수치 지표 계산, 시각화 생성 & RULE-AN-xxx, RULE-VZ-xxx \\
9. 대시보드 자동 생성 & 규칙 기반 레이아웃 구성 & RULE-VZ-xxx \\
10. AI 애널리스트 패널 & 템플릿 해석 → AI 자연어 보조 & RULE-IN-xxx, RULE-PO-xxx \\
11. Skills 로그 확인 & 5개 유형 로그 열람, 규칙 ID 추적 & 전체 RULE-xxx \\
\end{longtable}

\begin{center}\rule{0.5\linewidth}{0.5pt}\end{center}

\section{핵심 제품 원리: User = Data + Goal}\label{uxd575uxc2ec-uxc81cuxd488-uxc6d0uxb9ac-user-data-goal}

LENS는 사용자의 요청을 데이터셋 D와 분석 목적 G의 쌍으로 처리한다. 분석 실행 가능 여부는 D, G, 그리고 Skills.md 규칙의 함수다.

\subsection{목적-데이터 결합 매트릭스}\label{uxbaa9uxc801-uxb370uxc774uxd130-uxacb0uxd569-uxb9e4uxd2b8uxb9aduxc2a4}

\begin{longtable}[]{@{}
  >{\raggedright\arraybackslash}p{(\columnwidth - 6\tabcolsep) * \real{0.0882}}
  >{\raggedright\arraybackslash}p{(\columnwidth - 6\tabcolsep) * \real{0.3235}}
  >{\raggedright\arraybackslash}p{(\columnwidth - 6\tabcolsep) * \real{0.3235}}
  >{\raggedright\arraybackslash}p{(\columnwidth - 6\tabcolsep) * \real{0.2647}}@{}}
\toprule\noalign{}
\begin{minipage}[b]{\linewidth}\raggedright
\end{minipage} & \begin{minipage}[b]{\linewidth}\raggedright
데이터 충분
\end{minipage} & \begin{minipage}[b]{\linewidth}\raggedright
데이터 부분적
\end{minipage} & \begin{minipage}[b]{\linewidth}\raggedright
데이터 부족
\end{minipage} \\
\midrule\noalign{}
\endhead
\bottomrule\noalign{}
\endlastfoot
\textbf{목적 명확} & 전체 분석 실행 & 가능 분석 실행 + ``데이터 부족'' 패널 표시 & 분석 불가 항목과 필요 데이터 안내 \\
\textbf{목적 모호} & 분석 목적 후보 제시 → 사용자 선택 후 실행 & 분석 목적 후보 제시 → 선택 후 부분 분석 & 분석 목적 후보 제시 → 선택 후 필요 데이터 안내 \\
\end{longtable}

\subsection{설계 원칙}\label{uxc124uxacc4-uxc6d0uxce59}

\textbf{정책 1: 분석 목적 명확화 우선.}
분석 목적이 모호한 경우, 시스템은 분석을 즉시 시작하지 않는다. 성과 분석·위험 분석·비중 분석·자동 진단 중 선택지를 사용자에게 제시하며, 사용자 확인 없이 목적을 자동 확정하지 않는다.

\textbf{정책 2: 불가능한 분석은 사유를 표시하고, 가능한 분석은 계속 처리한다.}
데이터가 불충분한 경우에도 시스템은 오류를 출력하거나 분석을 전면 중단하지 않는다. 불가능한 분석 항목에는 사유와 필요 데이터를 표시하고, 가능한 분석 항목은 계속 처리하여 대시보드에 표시한다.

\textbf{정책 3: Skills.md 규칙이 분석 판단 기준을 결정한다.}
분석 가능 여부, 차트 유형 선택, 위험 경고 임계값은 모두 Skills.md에 사전 정의된 규칙에 따른다. AI가 임의로 계산 방식을 선택하거나 임계값을 설정하지 않는다.

\textbf{정책 4: 재현성 보장 범위를 수치 계산 영역으로 한정한다.}
동일 파일·동일 사용자 확정 매핑·동일 Skills.md 버전이 주어지면 수치 지표와 규칙 기반 시각화 구성이 재현된다. AI 애널리스트 패널의 자연어 표현은 실행마다 달라질 수 있으므로 재현성 보장 범위에서 제외한다.

\begin{center}\rule{0.5\linewidth}{0.5pt}\end{center}

\section{MVP 범위}\label{mvp-uxbc94uxc704}

본 기획서에서 P0는 공모전 데모에 반드시 포함할 핵심 기능, P1은 데모 이후 확장 또는 시간 여유 시 구현할 기능, P2는 장기 고도화 기능을 의미한다.

\subsection{P0 범위 (공모전 데모 핵심 경로)}\label{p0-uxbc94uxc704-uxacf5uxbaa8uxc804-uxb370uxbaa8-uxd575uxc2ec-uxacbduxb85c}

\textbf{지원 입력:}
- 단일 CSV/XLSX 파일 (10MB 이하)
- 5종목 1년치 일별 수익률 또는 일별 가격 데이터 (P0 데모 기준)
- 비중(weight) 컬럼: 존재 시 비중 분석 추가, 부재 시 비중 분석 패널 ``데이터 부족'' 상태 표시

\textbf{P0 분석 항목:}

\begin{longtable}[]{@{}lll@{}}
\toprule\noalign{}
분석 & 필요 컬럼 & 출력 \\
\midrule\noalign{}
\endhead
\bottomrule\noalign{}
\endlastfoot
누적 수익률 & return 또는 price & 선 차트 \\
구간 수익률 & return 또는 price & 수치 테이블 \\
최대낙폭(MDD) & return 또는 price & MDD 수치 + 드로다운 차트 \\
연환산 변동성 & return & 수치 테이블 \\
기본 비중 분석 & weight (존재 시) & 파이/바 차트 \\
데이터 충분성 결과 & --- & 항목별 판정 테이블 \\
AI 애널리스트 패널 & 위 지표 결과 & 규칙 기반 해석문 \\
Skills 로그 & 전 단계 처리 결과 & 5개 유형 탭 \\
\end{longtable}

\textbf{P0 명시적 제외 항목:}

\begin{longtable}[]{@{}
  >{\raggedright\arraybackslash}p{(\columnwidth - 4\tabcolsep) * \real{0.3333}}
  >{\raggedright\arraybackslash}p{(\columnwidth - 4\tabcolsep) * \real{0.3333}}
  >{\raggedright\arraybackslash}p{(\columnwidth - 4\tabcolsep) * \real{0.3333}}@{}}
\toprule\noalign{}
\begin{minipage}[b]{\linewidth}\raggedright
항목
\end{minipage} & \begin{minipage}[b]{\linewidth}\raggedright
분류
\end{minipage} & \begin{minipage}[b]{\linewidth}\raggedright
이유
\end{minipage} \\
\midrule\noalign{}
\endhead
\bottomrule\noalign{}
\endlastfoot
CAGR (핵심 지표로서) & P1 참고용 & 거래일수 $\geq 252$일 조건 미충족 시 표시 불가; 충족 시에도 경고 레이블 필수 참고용 표시만 허용 \\
샤프 지수 & P1 & 무위험 수익률 데이터 미포함 \\
집중도 분석 / HHI & P1 & Holdings 데이터 기반 별도 분석 \\
벤치마크 비교 & P1 & 벤치마크 수익률 컬럼 필요 \\
트리맵 & P1 & 비중 시각화 확장 기능 \\
Markdown 파일 다운로드 & P1 & P0는 화면 내 열람까지만 \\
PDF/PNG 내보내기 & P2 & P1 이후 확장 \\
다중 파일 결합 & P1 & 단일 파일 분석 우선 \\
거래내역 기반 성과 계산 & P1+ & 현금 흐름·수수료 처리 복잡도 \\
실시간 시세·자동매매 & 영구 제외 & 서비스 범위 외 \\
\end{longtable}

\subsection{P1 범위}\label{p1-uxbc94uxc704}

\begin{itemize}
\tightlist
\item
  Holdings/비중 독립 분석: 보유 종목 비중 기반 할당 테이블 + 파이/바 차트
\item
  벤치마크 비교: 파일 내 벤치마크 컬럼 존재 시 누적 수익률 비교 + 초과 수익 (트래킹 에러·정보 비율은 P2)
\item
  거래 내역 인식: 거래 목록 표시 + 기본 요약 (성과·MDD·변동성 역산은 제외)
\item
  Markdown 내보내기
\item
  다중 파일 업로드
\item
  CAGR: 거래일수 $\geq 252$일 충족 시 경고 레이블 포함 참고용 표시
\end{itemize}

\subsection{P2 범위}\label{p2-uxbc94uxc704}

\begin{itemize}
\tightlist
\item
  PDF/PNG 내보내기
\item
  고급 비중 시각화: 트리맵, 섹터 드릴다운, 집중도 지수(HHI), Top-N 비중
\item
  고급 벤치마크 지표: 트래킹 에러, 정보 비율, 알파·베타
\item
  대시보드 레이아웃 수동 편집
\end{itemize}

\begin{center}\rule{0.5\linewidth}{0.5pt}\end{center}

\section{P0 대표 데모 시나리오}\label{p0-uxb300uxd45c-uxb370uxbaa8-uxc2dcuxb098uxb9acuxc624}

\textbf{시나리오:} 포트폴리오 분석 담당자(페르소나 A)가 5종목 1년치 일별 수익률이 담긴 비정형 한국어 CSV 파일을 업로드해 월간 성과 리포트를 생성한다.

\begin{longtable}[]{@{}
  >{\raggedright\arraybackslash}p{(\columnwidth - 6\tabcolsep) * \real{0.1579}}
  >{\raggedright\arraybackslash}p{(\columnwidth - 6\tabcolsep) * \real{0.1579}}
  >{\raggedright\arraybackslash}p{(\columnwidth - 6\tabcolsep) * \real{0.2368}}
  >{\raggedright\arraybackslash}p{(\columnwidth - 6\tabcolsep) * \real{0.4474}}@{}}
\toprule\noalign{}
\begin{minipage}[b]{\linewidth}\raggedright
단계
\end{minipage} & \begin{minipage}[b]{\linewidth}\raggedright
화면
\end{minipage} & \begin{minipage}[b]{\linewidth}\raggedright
핵심 장면
\end{minipage} & \begin{minipage}[b]{\linewidth}\raggedright
심사위원 주목 포인트
\end{minipage} \\
\midrule\noalign{}
\endhead
\bottomrule\noalign{}
\endlastfoot
1 & SCR-01 & 컬럼명이 ``날짜'', ``종목코드'', ``수익률\_조정'', ``비중'' 등 비정형인 CSV 드롭 & 정의된 포트폴리오 데이터 유형의 CSV/XLSX라면 비정형 컬럼명도 처리 가능 \\
2 & SCR-02 & AI가 컬럼명을 date·ticker·return·weight로 자동 추론, 신뢰도 표시 & 비정형 입력의 자동 처리 \\
3 & SCR-02 & 사용자가 한 번 클릭으로 매핑 확정 & 수작업 없는 진입 \\
4 & SCR-03 & ``이 포트폴리오 괜찮은지 분석해줘'' 입력 & 자연어 목적 입력 \\
5 & SCR-03 & 시스템이 모호한 목적 감지 → 성과·위험·비중·자동 진단 후보 제시 & 목적 기반 라우팅 \\
6 & SCR-03 & 사용자가 ``자동 진단'' 선택 & 사용자 확인 후 확정 \\
7 & SCR-04 & 분석 항목별 판정 결과 표시 (누적 수익률·MDD·변동성 가능, 벤치마크 불가) & 데이터 충분성 검증 \\
8 & SCR-05 & 누적 수익률·MDD·드로다운·변동성·비중 대시보드 자동 생성 & Skills.md 규칙 기반 자동화 \\
9 & SCR-06 & AI 애널리스트 패널: 위험 레이블 + 해석문 + 데이터 한계 고지 & 규칙 기반 해석, 투자 권유 없음 \\
10 & SCR-07 & Skills 로그 팝업: 5개 유형 탭, 규칙 ID 확인 & 재현성 증거물 \\
\end{longtable}

\textbf{데모 첫 10초 장면:} 파일 드롭(2초) → 매핑 확인(3초) → 대시보드 생성(3초) → Skills 로그 팝업(2초). 이 흐름을 통해 심사위원은 Skills.md → 규칙 ID → 대시보드 카드의 연결을 직접 확인할 수 있다.

\begin{center}\rule{0.5\linewidth}{0.5pt}\end{center}

\section{Use Case 요약}\label{use-case-uxc694uxc57d}

\begin{longtable}[]{@{}
  >{\raggedright\arraybackslash}p{(\columnwidth - 10\tabcolsep) * \real{0.0784}}
  >{\raggedright\arraybackslash}p{(\columnwidth - 10\tabcolsep) * \real{0.2353}}
  >{\raggedright\arraybackslash}p{(\columnwidth - 10\tabcolsep) * \real{0.1569}}
  >{\raggedright\arraybackslash}p{(\columnwidth - 10\tabcolsep) * \real{0.2157}}
  >{\raggedright\arraybackslash}p{(\columnwidth - 10\tabcolsep) * \real{0.1176}}
  >{\raggedright\arraybackslash}p{(\columnwidth - 10\tabcolsep) * \real{0.1961}}@{}}
\toprule\noalign{}
\begin{minipage}[b]{\linewidth}\raggedright
UC
\end{minipage} & \begin{minipage}[b]{\linewidth}\raggedright
Use Case 명
\end{minipage} & \begin{minipage}[b]{\linewidth}\raggedright
우선순위
\end{minipage} & \begin{minipage}[b]{\linewidth}\raggedright
입력 데이터
\end{minipage} & \begin{minipage}[b]{\linewidth}\raggedright
출력
\end{minipage} & \begin{minipage}[b]{\linewidth}\raggedright
공모전 연결
\end{minipage} \\
\midrule\noalign{}
\endhead
\bottomrule\noalign{}
\endlastfoot
UC-01 & 일별 수익률 CSV 기반 성과·리스크 대시보드 생성 & P0 & 일별 수익률 CSV & 누적 수익률·MDD·변동성 대시보드 & 대시보드 자동 생성 (25점) \\
UC-02 & 모호한 분석 목적 명확화 & P0 & 모호한 자연어 목적 입력 & 분석 목적 후보 제시 → 사용자 선택 & 범용성 (25점) \\
UC-03 & 데이터 부족 시 가능·불가 분리 & P0 & 불완전한 CSV & 가능 분석 실행 + 불가 항목 ``데이터 부족'' 패널 & Skills.md 설계 (25점) \\
UC-04 & Skills 로그로 규칙 기반 분석 검증 & P0 & 대시보드 생성 결과 & 5개 유형 로그 열람, 규칙 ID 추적 & Skills.md 설계 (25점) \\
UC-05 & Holdings/비중 기반 포트폴리오 비중 분석 & P1 & 보유 종목 + 비중 파일 & 할당 테이블 + 파이/바 차트 & 범용성 (25점) \\
UC-06 & 벤치마크 데이터 포함 시 비교 분석 & P1 & 포트폴리오 + 벤치마크 컬럼 & 포트폴리오 vs 벤치마크 누적 수익률 + 초과 수익 & 범용성 (25점) \\
UC-07 & 거래 내역 인식 및 거래 요약 & P1 & 거래 내역 파일 & 거래 목록 + 기본 요약 (성과 계산 제외) & 범용성 (25점) \\
UC-08 & Markdown/PDF/PNG 리포트 내보내기 & P1/P2 & 대시보드 결과 & Markdown(P1), PDF/PNG(P2) & 실용성 (10점) \\
\end{longtable}

\begin{center}\rule{0.5\linewidth}{0.5pt}\end{center}

\section{User Flow 요약}\label{user-flow-uxc694uxc57d}

\subsection{주요 흐름 유형}\label{uxc8fcuxc694-uxd750uxb984-uxc720uxd615}

\begin{longtable}[]{@{}
  >{\raggedright\arraybackslash}p{(\columnwidth - 4\tabcolsep) * \real{0.3333}}
  >{\raggedright\arraybackslash}p{(\columnwidth - 4\tabcolsep) * \real{0.3333}}
  >{\raggedright\arraybackslash}p{(\columnwidth - 4\tabcolsep) * \real{0.3333}}@{}}
\toprule\noalign{}
\begin{minipage}[b]{\linewidth}\raggedright
흐름
\end{minipage} & \begin{minipage}[b]{\linewidth}\raggedright
조건
\end{minipage} & \begin{minipage}[b]{\linewidth}\raggedright
결과
\end{minipage} \\
\midrule\noalign{}
\endhead
\bottomrule\noalign{}
\endlastfoot
흐름 A: 정상 P0 경로 & 컬럼 고신뢰도, 목적 명확, 데이터 충분 & 전체 P0 대시보드 생성 \\
흐름 B: 목적 모호성 명확화 & 모호한 자연어 입력 & 후보 제시 → 선택 → 정상 경로 합류 \\
흐름 C: 데이터 불충분 & 필요 컬럼 부재 & 가능 분석 실행 + ``데이터 부족'' 패널 \\
흐름 D: 컬럼 매핑 불확실 & 중간/저신뢰도 컬럼 존재 & 팝업/드롭다운 확인 → 확정 후 진행 \\
흐름 E: Skills 로그 열람 & 대시보드 생성 후 & 5개 유형 탭 열람, 규칙 ID 확인 \\
\end{longtable}


\subsection{P0 핵심 흐름 다이어그램}\label{p0-핵심-흐름-다이어그램}

아래 도식은 P0 데모에서 사용자가 파일을 업로드한 뒤 대시보드와 Skills 로그를 확인하기까지의 핵심 흐름을 요약한 것이다.

\begin{figure}[htbp]
\centering
\resizebox{0.96\textwidth}{!}{%
\begin{tikzpicture}[
  node distance=7mm and 9mm,
  process/.style={rectangle, rounded corners, draw=black!70, fill=gray!8, align=center, minimum width=32mm, minimum height=8mm, font=\small},
  decision/.style={diamond, aspect=2.2, draw=black!70, fill=gray!5, align=center, inner sep=1pt, font=\small},
  arrow/.style={-{Latex[length=2mm]}, thick}
]
\node[process] (upload) {파일 업로드\\SCR-01};
\node[process, right=of upload] (preview) {데이터 미리보기\\컬럼 매핑\\SCR-02};
\node[decision, right=of preview] (conf) {컬럼\\신뢰도};
\node[process, above right=5mm and 10mm of conf] (auto) {고신뢰도\\일괄 확정};
\node[process, below right=5mm and 10mm of conf] (manual) {중/저신뢰도\\개별 확인};
\node[process, right=25mm of conf] (goal) {분석 목적 입력\\SCR-03};
\node[decision, right=of goal] (goalcheck) {목적\\명확?};
\node[process, above right=5mm and 9mm of goalcheck] (clarify) {목적 후보 제시\\사용자 선택};
\node[process, right=25mm of goalcheck] (suff) {데이터 충분성 검사\\SCR-04};
\node[decision, right=of suff] (datacheck) {데이터\\충분?};
\node[process, above right=5mm and 9mm of datacheck] (full) {전체 분석 실행};
\node[process, below right=5mm and 9mm of datacheck] (partial) {가능 분석 실행\\부족 항목 표시};
\node[process, right=26mm of datacheck] (dash) {대시보드 메인\\SCR-05};
\node[process, above right=5mm and 10mm of dash] (ai) {AI 애널리스트 패널\\SCR-06};
\node[process, below right=5mm and 10mm of dash] (log) {Skills 로그 패널\\SCR-07};

\draw[arrow] (upload) -- (preview);
\draw[arrow] (preview) -- (conf);
\draw[arrow] (conf) -- node[above, font=\scriptsize] {고} (auto);
\draw[arrow] (conf) -- node[below, font=\scriptsize] {중/저} (manual);
\draw[arrow] (auto) -- (goal);
\draw[arrow] (manual) -- (goal);
\draw[arrow] (goal) -- (goalcheck);
\draw[arrow] (goalcheck) -- node[above, font=\scriptsize] {모호} (clarify);
\draw[arrow] (clarify) -- (suff);
\draw[arrow] (goalcheck) -- node[below, font=\scriptsize] {명확} (suff);
\draw[arrow] (suff) -- (datacheck);
\draw[arrow] (datacheck) -- node[above, font=\scriptsize] {충분} (full);
\draw[arrow] (datacheck) -- node[below, font=\scriptsize] {부분} (partial);
\draw[arrow] (full) -- (dash);
\draw[arrow] (partial) -- (dash);
\draw[arrow] (dash) -- (ai);
\draw[arrow] (dash) -- (log);
\end{tikzpicture}%
}
\caption{P0 핵심 사용자 흐름}
\end{figure}

\subsection{P1/P2 확장 경계}\label{p1p2-uxd655uxc7a5-uxacbduxacc4}

대시보드 화면(SCR-05)에서 P1/P2 기능을 요청하면 SCR-08(P1/P2 확장 안내)로 이동한다. P0 분석 결과는 유지되며 확장 기능은 별도 안내만 표시된다.

\begin{center}\rule{0.5\linewidth}{0.5pt}\end{center}

\section{화면 설계 요약}\label{uxd654uxba74-uxc124uxacc4-uxc694uxc57d}


\subsection{전체 레이아웃}\label{전체-레이아웃}

\begin{table}[htbp]
\centering
\small
\begin{tabularx}{\textwidth}{|Y|Y|}
\hline
\multicolumn{2}{|l|}{\textbf{상태 바} \quad 파일명 \; | \; 매핑 상태 \; | \; 분석 목적 \; | \; 단계 진행 표시} \\
\hline
\textbf{대시보드 캔버스(좌측)} & \textbf{AI 애널리스트 패널(우측)} \\
\hline
\begin{itemize}[leftmargin=*, itemsep=1pt]
\item 누적 수익률 선 차트
\item 구간 수익률 테이블
\item MDD 및 드로다운 차트
\item 연환산 변동성 테이블
\item 기본 비중 차트(조건부)
\end{itemize}
&
\begin{itemize}[leftmargin=*, itemsep=1pt]
\item Skills.md 템플릿 해석문
\item 데이터 한계 고지
\item 위험 판정 레이블
\item 불가능한 분석과 필요 데이터 안내
\end{itemize}
\\
\hline
\multicolumn{2}{|l|}{\textbf{Skills 로그 슬라이드오버} \quad Column Mapping \; | \; Sufficiency \; | \; Metric \; | \; Visualization \; | \; Insight} \\
\hline
\end{tabularx}
\caption{LENS 메인 화면 구조}
\end{table}

좌측 영역은 수치와 차트를 표시하고, 우측 영역은 해석과 데이터 한계를 제공한다. 각 대시보드 카드 하단에는 참조 규칙 ID를 표시한다.

\subsection{화면 목록}\label{uxd654uxba74-uxbaa9uxb85d}

\begin{longtable}[]{@{}
  >{\raggedright\arraybackslash}p{(\columnwidth - 6\tabcolsep) * \real{0.2500}}
  >{\raggedright\arraybackslash}p{(\columnwidth - 6\tabcolsep) * \real{0.2222}}
  >{\raggedright\arraybackslash}p{(\columnwidth - 6\tabcolsep) * \real{0.1667}}
  >{\raggedright\arraybackslash}p{(\columnwidth - 6\tabcolsep) * \real{0.3611}}@{}}
\toprule\noalign{}
\begin{minipage}[b]{\linewidth}\raggedright
화면 ID
\end{minipage} & \begin{minipage}[b]{\linewidth}\raggedright
화면명
\end{minipage} & \begin{minipage}[b]{\linewidth}\raggedright
목적
\end{minipage} & \begin{minipage}[b]{\linewidth}\raggedright
핵심 구성 요소
\end{minipage} \\
\midrule\noalign{}
\endhead
\bottomrule\noalign{}
\endlastfoot
SCR-01 & 파일 업로드 & 파일 수신 및 형식 검증 & 드래그앤드롭 영역, 지원 형식 안내 \\
SCR-02 & 데이터 미리보기 \& 컬럼 매핑 & 매핑 확인 및 확정 & 미리보기 테이블, 컬럼-범주-신뢰도 매핑 표 \\
SCR-03 & 분석 목적 입력 & 목적 확정 및 모호성 해소 & 자연어 입력창, 목적 후보 카드 \\
SCR-04 & 데이터 충분성 결과 & 분석 가능 여부 명시 & 분석 항목별 판정 테이블 \\
SCR-05 & 대시보드 메인 & P0 분석 결과 시각화 & 차트 카드, 규칙 ID 배지 \\
SCR-06 & AI 애널리스트 패널 & 규칙 기반 해석문 표시 & 템플릿 해석문, 위험 레이블, 데이터 한계 \\
SCR-07 & Skills 로그 패널 & 규칙-결과 추적 & 5개 유형 탭, 필드 상세 열람 \\
SCR-08 & P1/P2 확장 안내 & P0 경계 외 기능 안내 & 기능 설명, 필요 추가 데이터 안내 \\
\end{longtable}

\begin{center}\rule{0.5\linewidth}{0.5pt}\end{center}

\section{기능 정의 요약}\label{uxae30uxb2a5-uxc815uxc758-uxc694uxc57d}

\subsection{P0 기능 목록}\label{p0-uxae30uxb2a5-uxbaa9uxb85d}

\begin{longtable}[]{@{}
  >{\raggedright\arraybackslash}p{(\columnwidth - 8\tabcolsep) * \real{0.2683}}
  >{\raggedright\arraybackslash}p{(\columnwidth - 8\tabcolsep) * \real{0.1463}}
  >{\raggedright\arraybackslash}p{(\columnwidth - 8\tabcolsep) * \real{0.2195}}
  >{\raggedright\arraybackslash}p{(\columnwidth - 8\tabcolsep) * \real{0.1463}}
  >{\raggedright\arraybackslash}p{(\columnwidth - 8\tabcolsep) * \real{0.2195}}@{}}
\toprule\noalign{}
\begin{minipage}[b]{\linewidth}\raggedright
Feature ID
\end{minipage} & \begin{minipage}[b]{\linewidth}\raggedright
기능
\end{minipage} & \begin{minipage}[b]{\linewidth}\raggedright
연관 화면
\end{minipage} & \begin{minipage}[b]{\linewidth}\raggedright
설명
\end{minipage} & \begin{minipage}[b]{\linewidth}\raggedright
수용 기준
\end{minipage} \\
\midrule\noalign{}
\endhead
\bottomrule\noalign{}
\endlastfoot
F-UP-01 & 단일 파일 업로드 & SCR-01 & CSV/XLSX 수신, 형식·인코딩·크기 검증 & 10MB 이하 CSV/XLSX → 미리보기 표시; 오류 시 원인 메시지 \\
F-CM-01 & AI 컬럼 매핑 추론 & SCR-02 & 컬럼명 → 의미 범주 추론 + 신뢰도 & 고/중/저 3단계 신뢰도 구분 표시 \\
F-CM-02 & 신뢰도별 UI 분기 & SCR-02 & 고: 일괄 확인, 중: 개별 팝업, 저: 수동 드롭다운 & 각 구간별 지정 UI 동작 \\
F-CM-03 & 매핑 확정 & SCR-02 & 사용자 확인 후 확정, Column Mapping Log 기록 & 확정 시 로그에 원본명·범주·신뢰도·확정 시각 포함 \\
F-GC-01 & 분석 목적 입력 & SCR-03 & 자연어 입력 또는 후보 카드 선택 & 목적 범주 확정 후 충분성 검사 진행 \\
F-GC-02 & 목적 모호성 명확화 & SCR-03 & 모호한 입력 시 후보 제시 → 사용자 확인 & 사용자 확인 없이 목적 일방 확정 금지 \\
F-DS-01 & 데이터 충분성 검사 & SCR-04 & 분석 항목별 필요 컬럼 존재 여부 판정 & 가능/제한/불가 3단계, 불가 항목에 필요 데이터 명시 \\
F-DS-02 & 부분 분석 실행 & SCR-04→05 & 불가 항목 ``데이터 부족'' 패널, 가능 항목만 실행 & 침묵·오류 금지; 사유와 함께 표시 \\
F-AN-01 & 누적 수익률 계산 & SCR-05 & Skills.md 공식 기반 계산 & 동일 입력·매핑·규칙에서 결과 재현 가능 \\
F-AN-02 & 구간 수익률 계산 & SCR-05 & 일간·월간·전체 구간 & 결정적 출력 \\
F-AN-03 & MDD 및 드로다운 계산 & SCR-05 & 최대낙폭 수치 + 드로다운 시계열 & MDD 수치와 차트 모두 표시 \\
F-AN-04 & 연환산 변동성 계산 & SCR-05 & $\operatorname{std}(r_t)\times\sqrt{252}$ & 계산 공식 Metric Log에 기록 \\
F-AN-05 & 기본 비중 분석 & SCR-05 & weight 컬럼 존재 시 할당 테이블 + 파이/바 차트 & 컬럼 부재 시 ``데이터 부족'' 패널 \\
F-DB-01 & 대시보드 자동 생성 & SCR-05 & Skills.md 시각화 규칙 기반 차트 유형·레이아웃 자동 선택 & 카드 하단 참조 규칙 ID 표시 \\
F-AI-01 & AI 애널리스트 패널 & SCR-06 & Skills.md 템플릿 → AI 자연어 보조 (금지 표현 검증 후) & 투자 권유·미래 전망·인과 단정 금지 \\
F-SL-01 & Skills 로그 & SCR-07 & 5개 유형 로그 기록 및 탭 열람, 규칙 ID 양방향 연결 & 각 분석 단계마다 해당 유형 로그 기록 \\
\end{longtable}

\subsection{P1/P2 기능 경계}\label{p1p2-uxae30uxb2a5-uxacbduxacc4}

\begin{longtable}[]{@{}
  >{\raggedright\arraybackslash}p{(\columnwidth - 4\tabcolsep) * \real{0.2857}}
  >{\raggedright\arraybackslash}p{(\columnwidth - 4\tabcolsep) * \real{0.2857}}
  >{\raggedright\arraybackslash}p{(\columnwidth - 4\tabcolsep) * \real{0.4286}}@{}}
\toprule\noalign{}
\begin{minipage}[b]{\linewidth}\raggedright
기능
\end{minipage} & \begin{minipage}[b]{\linewidth}\raggedright
분류
\end{minipage} & \begin{minipage}[b]{\linewidth}\raggedright
진입 조건
\end{minipage} \\
\midrule\noalign{}
\endhead
\bottomrule\noalign{}
\endlastfoot
Holdings 비중 분석 (할당 테이블 + 파이/바 차트) & P1 & 보유 비중 파일 업로드 \\
벤치마크 비교 & P1 & 파일 내 벤치마크 컬럼 존재 \\
거래 내역 요약 (목록·요약만, 성과 계산 제외) & P1 & 거래 내역 파일 업로드 \\
CAGR 참고용 표시 & P1 & 거래일수 $\geq 252$일 + 경고 레이블 필수 \\
Markdown 내보내기 & P1 & P0 대시보드 완성 후 \\
다중 파일 업로드 & P1 & P0 이후 \\
PDF/PNG 내보내기 & P2 & P1 이후 \\
고급 비중 시각화 (트리맵, HHI, Top-N) & P2 & P1 Holdings 이후 \\
고급 벤치마크 지표 (트래킹 에러, 정보 비율) & P2 & P1 벤치마크 이후 \\
\end{longtable}

\begin{center}\rule{0.5\linewidth}{0.5pt}\end{center}

\section{Skills.md 설계}\label{skills.md-uxc124uxacc4}

\subsection{Skills.md의 역할}\label{skills.mduxc758-uxc5eduxd560}

본 기획서에서 Skills.md는 분석 처리 흐름 전체의 판단 기준을 사전 정의하는 \textbf{규칙 명세서}로 설계한다. AI는 이 규칙 범위 안에서만 작동한다.

Skills.md가 사전 정의하는 규칙 영역:
- 비정형 컬럼명을 분석 의미 단위로 추론하는 기준 (컬럼 매핑 규칙)
- 모호한 목적 입력을 분석 가능 범주로 변환하는 기준 (목적 명확화 규칙)
- 분석 항목별 필요 컬럼·데이터 기간 조건 (데이터 충분성 규칙)
- 각 지표의 계산 공식과 거래일수 기준 (지표 계산 규칙)
- 데이터 유형별 차트 유형 결정 로직 (시각화 선택 규칙)
- 위험 판정 임계값과 해석문 템플릿 (인사이트 템플릿 규칙)
- AI 출력에서 허용되지 않는 표현 유형 (금지 표현 규칙)
- 결정적 출력 보장 범위와 비결정 구간 (재현성 계약)

\subsection{Skills.md 구조 (제안 개요)}\label{skills.md-uxad6cuxc870-uxc81cuxc548-uxac1cuxc694}

\begin{verbatim}
Skills.md
├── metadata          # 버전, 서비스명, 최종 수정일
├── supported_input_types
│   ├── daily_returns     # 일별 수익률
│   ├── daily_prices      # 일별 가격
│   ├── holdings_weights  # 보유 비중 (P1)
│   └── transaction_log   # 거래 내역 (P1: 목록·요약만)
├── column_mapping_rules       # RULE-CM-xxx
├── goal_clarification_rules   # RULE-GC-xxx
├── data_sufficiency_rules     # RULE-DS-xxx
├── metric_rules               # RULE-AN-xxx
├── visualization_rules        # RULE-VZ-xxx
├── insight_templates          # RULE-IN-xxx
├── prohibited_outputs         # RULE-PO-xxx
├── skills_log_schema          # 5개 유형 필드 명세
└── reproducibility_contract   # 결정적/비결정적 경계 명시
\end{verbatim}


\subsection{규칙 예시}\label{규칙-예시}

\begin{description}[leftmargin=0pt, style=nextline, itemsep=0.6em]
\item[\textbf{RULE-CM-001: 수익률 컬럼 추론}]
인식 패턴은 ``수익률'', ``returns'', ``pnl'', ``ret'', ``일별수익률''로 둔다. 신뢰도 기준은 고($\geq 0.8$), 중($0.5 \leq s < 0.8$), 저($s<0.5$)로 구분한다. 고신뢰도는 자동 확정 후보로 제시하고, 중·저신뢰도는 사용자 확인 또는 수동 선택을 요구한다.

\item[\textbf{RULE-DS-001: MDD 충분성 조건}]
MDD 계산에는 \texttt{return} 또는 \texttt{price} 컬럼과 최소 30거래일 이상의 관측치가 필요하다. 조건을 충족하지 못하면 해당 분석을 불가로 판정하고, \texttt{excluded\_analysis\_reason}에 ``return/price 컬럼 없음'' 또는 ``관측 기간 부족''을 기록한다.

\item[\textbf{RULE-AN-001: 누적 수익률}]
\[
R_{\mathrm{cum}} = \prod_{t=1}^{T}(1+r_t)-1
\]
동일 입력, 동일 매핑, 동일 Skills.md 버전이 주어지면 계산 결과는 재현된다.

\item[\textbf{RULE-AN-002: 최대낙폭(MDD)}]
\[
\mathrm{MDD}=\min_t \left(\frac{V_t-\max_{0\leq s\leq t}V_s}{\max_{0\leq s\leq t}V_s}\right)
\]
MDD 값과 드로다운 시계열은 Metric Calculation Log에 함께 기록한다.

\item[\textbf{RULE-AN-003: 연환산 변동성}]
\[
\sigma_{\mathrm{ann}} = \operatorname{std}(r_t)\times\sqrt{252}
\]
CAGR은 거래일수 252일 이상일 때만 참고 지표로 표시하며, P0 핵심 지표에는 포함하지 않는다.

\item[\textbf{RULE-VZ-001: 시계열 시각화 선택}]
시계열 성과 추이는 선 차트로 표시한다. Visualization Rule Log에는 \texttt{chart\_selection\_reason} 필드로 ``시계열 추이 표현에 선 차트가 적합''을 기록한다.

\item[\textbf{RULE-IN-001: MDD 위험 판정 템플릿}]
MDD가 $-20\%$보다 낮으면 ``주의'', $-20\%$ 이상 $-10\%$ 미만이면 ``경계'', $-10\%$ 이상이면 ``정상'' 레이블을 부여한다. AI는 템플릿 결과를 자연어로 다듬는 역할만 수행한다.

\item[\textbf{RULE-PO-001: 금지 표현}]
투자 권유, 미래 수익률 전망, 근거 없는 인과 관계 단정, 업로드 데이터 외 시장 설명은 출력하지 않는다. 위반 문장이 발견되면 해당 문장을 제거하고, 재출력에 실패하면 사전 정의된 템플릿 문장으로 대체한다.
\end{description}

\begin{center}\rule{0.5\linewidth}{0.5pt}\end{center}

\section{AI 오케스트레이션 설계}\label{ai-uxc624uxcf00uxc2a4uxd2b8uxb808uxc774uxc158-uxc124uxacc4}

본 섹션은 실제 배포 마이크로서비스 명세가 아닌 기능 모듈 설계다. 각 모듈은 독립적 역할을 가지며, 이전 단계의 확정 출력만을 입력으로 사용한다.

\begin{longtable}[]{@{}
  >{\raggedright\arraybackslash}p{(\columnwidth - 8\tabcolsep) * \real{0.1667}}
  >{\raggedright\arraybackslash}p{(\columnwidth - 8\tabcolsep) * \real{0.1667}}
  >{\raggedright\arraybackslash}p{(\columnwidth - 8\tabcolsep) * \real{0.1667}}
  >{\raggedright\arraybackslash}p{(\columnwidth - 8\tabcolsep) * \real{0.2500}}
  >{\raggedright\arraybackslash}p{(\columnwidth - 8\tabcolsep) * \real{0.2500}}@{}}
\toprule\noalign{}
\begin{minipage}[b]{\linewidth}\raggedright
모듈
\end{minipage} & \begin{minipage}[b]{\linewidth}\raggedright
입력
\end{minipage} & \begin{minipage}[b]{\linewidth}\raggedright
출력
\end{minipage} & \begin{minipage}[b]{\linewidth}\raggedright
규칙 의존
\end{minipage} & \begin{minipage}[b]{\linewidth}\raggedright
실패 처리
\end{minipage} \\
\midrule\noalign{}
\endhead
\bottomrule\noalign{}
\endlastfoot
Data Profiler & 업로드 파일 & 컬럼 목록, 행 수, 인코딩 요약 & 지원 형식 목록 & 형식 오류 → 오류 메시지 + 조치 안내 \\
Column Mapper & 컬럼 목록 & 컬럼-범주 매핑 후보 + 신뢰도 & RULE-CM-xxx & 전체 저신뢰도 → 수동 입력 유도 \\
Goal Clarifier & 사용자 목적 & 확정된 분석 목적 범주 & RULE-GC-xxx & 모호 → 후보 카드 제시 \\
Sufficiency Checker & 확정 매핑 + 목적 & 항목별 판정 + \texttt{excluded\_analysis\_reason} & RULE-DS-xxx & 전체 불가 → 파일 유형 변경 안내 \\
Metric Engine & 확정 매핑 + 가능 분석 목록 & 지표 계산 결과 + rule\_input/output & RULE-AN-xxx & 계산 오류 → 해당 패널만 ``오류'' 표시 \\
Dashboard Composer & 계산 결과 + 분석 유형 & 카드 레이아웃 + 차트 유형 & RULE-VZ-xxx & 규칙 미매칭 → 기본 선 차트 fallback \\
Analyst Panel Generator & 계산 결과 + 인사이트 템플릿 & 템플릿 해석문 → AI 자연어 보조 & RULE-IN-xxx & AI 오류 → 템플릿 문장만 표시 \\
Skills Logger & 각 단계 처리 결과 & 5개 유형 로그 기록 & 전체 RULE & 로그 실패 → 분석 유지, 경고 표시 \\
Critic/Guardrail Checker & AI 해석문 초안 & 금지 표현 검증 통과 + 정제 해석문 & RULE-PO-xxx & 위반 → 해당 문장 제거 재출력 \\
\end{longtable}

\textbf{사람 개입 지점 (Human-in-the-loop):}
1. 컬럼 매핑 확인 (SCR-02): 사용자 확인 없이 자동 확정 금지
2. 분석 목적 명확화 (SCR-03): 모호한 목적의 일방 확정 금지
3. 가능한 분석으로 진행 확정 (SCR-04): 사용자 확인 후 진행
4. Skills 로그 열람 (SCR-07): 선택적, 항상 접근 가능

\begin{center}\rule{0.5\linewidth}{0.5pt}\end{center}

\section{데이터 충분성 판단 정책}\label{uxb370uxc774uxd130-uxcda9uxbd84uxc131-uxd310uxb2e8-uxc815uxcc45}

시스템은 다음 순서로 충분성을 판단한다.

\begin{enumerate}
\def\labelenumi{\arabic{enumi}.}
\tightlist
\item
  \textbf{필요 컬럼 존재 여부:} 요청 분석에 필요한 컬럼이 파일 내에 있는가
\item
  \textbf{대체 컬럼 가능성:} 없는 경우, 가격 컬럼에서 수익률을 역산할 수 있는가
\item
  \textbf{부분 분석 가능성:} 대체도 불가한 경우, 일부 지표만 계산 가능한가
\item
  \textbf{필요 데이터 안내:} 부분 분석도 불가한 경우, 어떤 추가 데이터가 필요한지 명시
\end{enumerate}

\textbf{판단 예시:}

\begin{longtable}[]{@{}
  >{\raggedright\arraybackslash}p{(\columnwidth - 2\tabcolsep) * \real{0.3750}}
  >{\raggedright\arraybackslash}p{(\columnwidth - 2\tabcolsep) * \real{0.6250}}@{}}
\toprule\noalign{}
\begin{minipage}[b]{\linewidth}\raggedright
상황
\end{minipage} & \begin{minipage}[b]{\linewidth}\raggedright
시스템 반응
\end{minipage} \\
\midrule\noalign{}
\endhead
\bottomrule\noalign{}
\endlastfoot
return 컬럼 없고 price 컬럼 있음 & 수익률 역산 후 분석 진행, 역산 사실 로그 기록 \\
벤치마크 비교 요청 + 벤치마크 컬럼 없음 & 포트폴리오 분석만 실행, 벤치마크 비교는 ``데이터 부족'' 패널 \\
weight 컬럼 없음 & 비중 분석 패널 ``데이터 부족'' 표시, 다른 분석은 정상 실행 \\
return + price 모두 없음 & 성과·위험 분석 불가 안내 + 필요 컬럼 명시 \\
거래 내역 파일만 존재 & 거래 목록 + 기본 요약만 (수익률·MDD·변동성 역산은 P1+) \\
\end{longtable}

불가능한 분석은 사유와 필요 데이터를 ``데이터 부족'' 패널에 명시하고, 가능한 분석은 계속 처리한다. 오류 출력이나 공백 화면으로 처리하지 않는다.

\begin{center}\rule{0.5\linewidth}{0.5pt}\end{center}

\section{재현성 및 감사 가능성}\label{uxc7acuxd604uxc131-uxbc0f-uxac10uxc0ac-uxac00uxb2a5uxc131}

\subsection{결정적 처리 vs AI 보조 처리 경계}\label{uxacb0uxc815uxc801-uxcc98uxb9ac-vs-ai-uxbcf4uxc870-uxcc98uxb9ac-uxacbduxacc4}

\begin{longtable}[]{@{}
  >{\raggedright\arraybackslash}p{(\columnwidth - 4\tabcolsep) * \real{0.3103}}
  >{\raggedright\arraybackslash}p{(\columnwidth - 4\tabcolsep) * \real{0.3793}}
  >{\raggedright\arraybackslash}p{(\columnwidth - 4\tabcolsep) * \real{0.3103}}@{}}
\toprule\noalign{}
\begin{minipage}[b]{\linewidth}\raggedright
처리 단계
\end{minipage} & \begin{minipage}[b]{\linewidth}\raggedright
결정적 여부
\end{minipage} & \begin{minipage}[b]{\linewidth}\raggedright
제어 정책
\end{minipage} \\
\midrule\noalign{}
\endhead
\bottomrule\noalign{}
\endlastfoot
파일 파싱 & 결정적 & 지원 형식 목록 기반 \\
컬럼 매핑 추론 (초기) & 비결정적 & 신뢰도 표시, 사용자 확인 필수 \\
컬럼 매핑 (사용자 확정 후) & 결정적 & Column Mapping Log에 확정 시각 기록 \\
데이터 충분성 판정 & 결정적 & RULE-DS-xxx 기준 \\
지표 계산 & 결정적 & RULE-AN-xxx 공식 고정 \\
시각화 선택 & 결정적 & RULE-VZ-xxx 기준 \\
AI 애널리스트 해석문 & 비결정적 & 금지 표현 검증(RULE-PO-001) 후 출력 \\
Skills 로그 기록 & 결정적 & 수정 불가 형태로 자동 기록 \\
\end{longtable}

\subsection{재현성 계약}\label{uxc7acuxd604uxc131-uxacc4uxc57d}

동일 파일, 동일 사용자 확정 매핑, 동일 Skills.md 버전이 주어지면 수치 지표와 규칙 기반 시각화 구성이 재현된다. AI 애널리스트 패널의 자연어 표현은 실행마다 달라질 수 있으므로 재현성 보장 범위에서 제외한다.

재현성이 보장되는 항목: 누적 수익률, 구간 수익률, MDD, 드로다운 시계열, 연환산 변동성, 시각화 차트 유형 선택, 대시보드 레이아웃 구성 규칙, Skills 로그

재현성 범위 외 항목: AI 애널리스트 패널 자연어 서술, 초기 컬럼 매핑 추론 결과 (사용자 확정 전)

\begin{center}\rule{0.5\linewidth}{0.5pt}\end{center}

\section{Skills 로그 설계}\label{skills-uxb85cuxadf8-uxc124uxacc4}

Skills 로그는 Skills.md 규칙이 실제 대시보드 생성에 적용되었음을 기록하는 감사 항목이다. 대시보드 카드 하단에 표시된 규칙 ID를 클릭하면 해당 로그 항목으로 이동하는 양방향 연결 구조를 제공한다.

\begin{longtable}[]{@{}
  >{\raggedright\arraybackslash}p{(\columnwidth - 6\tabcolsep) * \real{0.3000}}
  >{\raggedright\arraybackslash}p{(\columnwidth - 6\tabcolsep) * \real{0.2000}}
  >{\raggedright\arraybackslash}p{(\columnwidth - 6\tabcolsep) * \real{0.3000}}
  >{\raggedright\arraybackslash}p{(\columnwidth - 6\tabcolsep) * \real{0.2000}}@{}}
\toprule\noalign{}
\begin{minipage}[b]{\linewidth}\raggedright
로그 유형
\end{minipage} & \begin{minipage}[b]{\linewidth}\raggedright
트리거
\end{minipage} & \begin{minipage}[b]{\linewidth}\raggedright
필수 필드
\end{minipage} & \begin{minipage}[b]{\linewidth}\raggedright
역할
\end{minipage} \\
\midrule\noalign{}
\endhead
\bottomrule\noalign{}
\endlastfoot
Column Mapping Log & 매핑 확정 (SCR-02) & 원본 컬럼명, 추론 범주, 신뢰도, 사용자 확인 여부, 확정 시각 & 컬럼 해석 투명성 \\
Sufficiency Check Log & 충분성 판정 (SCR-04) & 분석 항목명, 판정 결과, 부족 데이터, 규칙 ID, \texttt{excluded\_analysis\_reason} & 분석 제외 근거 추적 \\
Metric Calculation Log & 지표 계산 완료 (SCR-05) & 지표명, \texttt{calculation\_formula}, \texttt{rule\_input}, \texttt{rule\_output}, 규칙 ID & 수치 재현성 근거 \\
Visualization Rule Log & 차트 생성 (SCR-05) & 분석 유형, 차트 유형, \texttt{chart\_selection\_reason}, 규칙 ID & 시각화 선택 근거 \\
Insight Rule Log & AI 패널 출력 (SCR-06) & 참조 지표, 임계값 규칙, 위험 판정 결과, AI 보조 문장 여부 & 해석 근거와 AI 개입 범위 \\
\end{longtable}

Metric Calculation Log의 \texttt{calculation\_formula}, \texttt{rule\_input}, \texttt{rule\_output} 필드를 통해 심사위원이 계산 과정을 직접 확인할 수 있다. 이 항목은 Skills.md 규칙이 실제 수치 산출에 적용되었음을 확인하는 근거가 된다.

\begin{center}\rule{0.5\linewidth}{0.5pt}\end{center}

\section{AI 애널리스트 패널 설계}\label{ai-uxc560uxb110uxb9acuxc2a4uxd2b8-uxd328uxb110-uxc124uxacc4}

\subsection{작동 원칙}\label{uxc791uxb3d9-uxc6d0uxce59}

1차: Skills.md RULE-IN-xxx 템플릿이 트리거 조건(임계값 초과 여부)을 판단하고 해석문 구조를 결정한다.
2차: AI가 템플릿 결과를 자연어로 polishing한다.
3차: Critic/Guardrail Checker가 금지 표현 검증을 수행하고 통과된 문장만 출력한다.

\subsection{금지 출력 유형}\label{uxae08uxc9c0-uxcd9cuxb825-uxc720uxd615}

\begin{itemize}
\tightlist
\item
  투자 권유 표현 (``매수'', ``매도'', ``보유 추천'')
\item
  미래 수익률 전망 (``향후 수익'', ``상승 예상'')
\item
  근거 없는 인과 관계 단정 (``A 때문에 B가 하락'')
\item
  업로드 데이터 외 외부 시장 설명 (데이터에 포함되지 않은 시장 상황 서술)
\end{itemize}

\subsection{패널 구성 항목}\label{uxd328uxb110-uxad6cuxc131-uxd56duxbaa9}

\begin{longtable}[]{@{}
  >{\raggedright\arraybackslash}p{(\columnwidth - 2\tabcolsep) * \real{0.5000}}
  >{\raggedright\arraybackslash}p{(\columnwidth - 2\tabcolsep) * \real{0.5000}}@{}}
\toprule\noalign{}
\begin{minipage}[b]{\linewidth}\raggedright
항목
\end{minipage} & \begin{minipage}[b]{\linewidth}\raggedright
내용
\end{minipage} \\
\midrule\noalign{}
\endhead
\bottomrule\noalign{}
\endlastfoot
면책 고지 & ``이 리포트는 투자 권유가 아니며, 업로드된 데이터의 분석 결과를 요약한 것이다'' \\
관찰된 주요 지표 & 누적 수익률, MDD, 연환산 변동성 수치 요약 \\
위험 레이블 & Skills.md RULE-IN-001 기준 위험 판정 결과 \\
데이터 한계 & 분석 기간, 데이터 수, 부재 컬럼 기반 분석 제한 사항 \\
불가능한 분석 안내 & Sufficiency Check 결과 기반 추가 필요 데이터 명시 \\
Skills 로그 참조 & ``이 해석은 RULE-IN-001, RULE-IN-002를 따릅니다'' \\
\end{longtable}

\begin{center}\rule{0.5\linewidth}{0.5pt}\end{center}

\section{공모전 평가항목별 대응 전략}\label{uxacf5uxbaa8uxc804-uxd3c9uxac00uxd56duxbaa9uxbcc4-uxb300uxc751-uxc804uxb7b5}

\begin{longtable}[]{@{}
  >{\raggedright\arraybackslash}p{(\columnwidth - 6\tabcolsep) * \real{0.2045}}
  >{\raggedright\arraybackslash}p{(\columnwidth - 6\tabcolsep) * \real{0.1136}}
  >{\raggedright\arraybackslash}p{(\columnwidth - 6\tabcolsep) * \real{0.3636}}
  >{\raggedright\arraybackslash}p{(\columnwidth - 6\tabcolsep) * \real{0.3182}}@{}}
\toprule\noalign{}
\begin{minipage}[b]{\linewidth}\raggedright
평가항목
\end{minipage} & \begin{minipage}[b]{\linewidth}\raggedright
배점
\end{minipage} & \begin{minipage}[b]{\linewidth}\raggedright
LENS 대응
\end{minipage} & \begin{minipage}[b]{\linewidth}\raggedright
데모에서의 증거
\end{minipage} \\
\midrule\noalign{}
\endhead
\bottomrule\noalign{}
\endlastfoot
범용성 & 25 & 정의된 CSV/XLSX 데이터 유형 내에서 컬럼 구조가 달라도 처리; RULE-CM-xxx가 비정형 컬럼명 추론 & 비정형 한국어 컬럼명 CSV 업로드 후 자동 매핑 장면 \\
Skills.md 설계 & 25 & 8개 규칙 영역, 규칙 ID 체계, 예시 규칙 블록, Skills 로그 5개 유형이 명세서 수준으로 제출됨 & Skills 로그 팝업에서 규칙 ID와 계산 공식 직접 확인 \\
대시보드 자동 생성 & 25 & 9개 모듈 처리 흐름이 업로드부터 대시보드까지 수작업 없이 자동 완주; RULE-VZ-xxx가 차트 유형 자동 선택 & 업로드 → 대시보드 P0 샘플 데이터 기준 30초 이내 완주 장면 \\
바이브코딩 활용 & 15 & Skills.md를 기준 문서로 삼아 AI 개발 도구가 화면·기능·분석 규칙 산출물을 일관되게 생성하고, 리뷰 에이전트가 P0 범위·금지 표현·재현성 조건을 검증하는 방식으로 바이브코딩을 활용한다 (AI 개발 도구 예시: Claude Code Subagents) & Skills.md 기준으로 생성된 코드·대시보드 연결 관계 및 리뷰 로그 제시 \\
실용성 및 창의성 & 10 & 데이터 충분성 검사 + ``침묵 대신 부분 분석'' + Skills 로그 가시화가 실무 활용 가능성을 높임 & 벤치마크 컬럼 부재 시 ``데이터 부족'' 패널 표시 장면 \\
\end{longtable}

\textbf{범용성 범위 정의:} 본 기획서에서의 범용성은 ``정의된 4개 데이터 유형(일별 수익률, 일별 가격, 보유 비중, 거래 내역) 내에서 컬럼 구조가 다양해도 처리할 수 있는 유연성''으로 한정한다. 임의 투자 데이터 전체에 대한 범용 처리는 본 범위에 포함되지 않는다.

\begin{center}\rule{0.5\linewidth}{0.5pt}\end{center}

\section{기존 서비스 대비 차별화}\label{uxae30uxc874-uxc11cuxbe44uxc2a4-uxb300uxbe44-uxcc28uxbcc4uxd654}

\begin{quote}
\textbf{면책 표기:} 이하 서술은 각 서비스의 공개된 제품 설명을 기반으로 한 기획 해석이며, 해당 서비스의 공식 결함 선언이 아니다.
\end{quote}

\begin{longtable}[]{@{}
  >{\raggedright\arraybackslash}p{(\columnwidth - 4\tabcolsep) * \real{0.1875}}
  >{\raggedright\arraybackslash}p{(\columnwidth - 4\tabcolsep) * \real{0.3542}}
  >{\raggedright\arraybackslash}p{(\columnwidth - 4\tabcolsep) * \real{0.4583}}@{}}
\toprule\noalign{}
\begin{minipage}[b]{\linewidth}\raggedright
비교 대상
\end{minipage} & \begin{minipage}[b]{\linewidth}\raggedright
공개 정보 기반 특징
\end{minipage} & \begin{minipage}[b]{\linewidth}\raggedright
LENS의 차별화 설계
\end{minipage} \\
\midrule\noalign{}
\endhead
\bottomrule\noalign{}
\endlastfoot
범용 AI 분석 도구 (Julius AI 등) & CSV 업로드 + 자연어 분석 지원. 금융 지표를 포함한 다양한 분석이 프롬프트 기반으로 가능 & Skills.md 기반 금융 포트폴리오 전용 규칙이 분석 흐름 전체를 제어하는 구조가 공개 사용 흐름에서 전면화되어 있지 않다. LENS는 이를 공모전 요구사항에 맞춰 전면 기능으로 설계한다 \\
범용 BI + AI (Power BI Copilot 등) & CSV 업로드와 자연어 대시보드 생성 지원. Microsoft 공식 문서는 Copilot 출력의 비결정성을 명시 & Skills.md가 금융 포트폴리오 전용 규칙을 사전 정의하고 Skills 로그가 규칙-결과 연결을 기록한다. 결정적 수치 계산 영역에서 재현성을 보장한다 \\
금융 특화 분석 서비스 (Portfolio Visualizer, Koyfin, Sharesight 등) & 티커 입력, 브로커 연동, 자체 포트폴리오 등록을 주된 데이터 진입 경로로 제시 (기획 해석) & 비정형 CSV/XLSX를 그대로 업로드해 AI 컬럼 매핑으로 진입하는 경로가 공개 사용 흐름의 전면에 제시되어 있지 않다. LENS는 이를 진입 경로 설계의 핵심으로 삼는다 \\
\end{longtable}

\textbf{LENS의 설계 차별점 (기획 해석):}

\begin{enumerate}
\def\labelenumi{\arabic{enumi}.}
\tightlist
\item
  \textbf{Data + Goal 동시 해석:} 파일과 분석 목적을 함께 입력받아, 충분성 판단 후 가능한 분석만 실행한다
\item
  \textbf{분석 전 충분성 검사:} 실행 전에 분석 항목별 가능 여부를 명시한다
\item
  \textbf{부분 분석 처리:} 불가능한 분석 항목에는 사유와 필요 데이터를 표시하고, 가능한 항목은 계속 처리한다
\item
  \textbf{Skills.md 규칙 사전 명세:} 계산 공식·임계값·차트 선택·인사이트 기준을 모두 사전에 문서화한다
\item
  \textbf{Skills 로그에 감사 근거 기록:} 규칙 ID·입출력·판정 근거가 모든 분석 단계에 기록된다
\end{enumerate}

\begin{center}\rule{0.5\linewidth}{0.5pt}\end{center}

\section{리스크 및 통제 정책}\label{uxb9acuxc2a4uxd06c-uxbc0f-uxd1b5uxc81c-uxc815uxcc45}

\begin{longtable}[]{@{}
  >{\raggedright\arraybackslash}p{(\columnwidth - 4\tabcolsep) * \real{0.3478}}
  >{\raggedright\arraybackslash}p{(\columnwidth - 4\tabcolsep) * \real{0.2609}}
  >{\raggedright\arraybackslash}p{(\columnwidth - 4\tabcolsep) * \real{0.3913}}@{}}
\toprule\noalign{}
\begin{minipage}[b]{\linewidth}\raggedright
리스크
\end{minipage} & \begin{minipage}[b]{\linewidth}\raggedright
영향
\end{minipage} & \begin{minipage}[b]{\linewidth}\raggedright
통제 정책
\end{minipage} \\
\midrule\noalign{}
\endhead
\bottomrule\noalign{}
\endlastfoot
P0 범위 과도한 확장 & 데모 완주 실패 & SCR-08(P1/P2 확장 안내)로 경계 명시; P0 분석 목록 고정 \\
AI 과장 서술 (투자 권유 등) & 신뢰 손상, 심사 감점 & RULE-PO-001 + Critic/Guardrail Checker 강제 적용 \\
목적 모호성 미처리 & 엉뚱한 분석 실행 & 목적 일방 확정 금지; 후보 카드 제시 필수 \\
데이터 불충분 무시 & 오류·침묵 출력 & Sufficiency Checker 필수 단계, ``데이터 부족'' 패널 항상 표시 \\
컬럼 매핑 오류 & 잘못된 분석 결과 & 신뢰도별 UI 분기, 사용자 확정 전 자동 분석 금지 \\
재현성 과도한 주장 & 비결정적 AI 영역과 충돌 & 재현성 계약 범위를 ``수치 계산 영역''으로 명시 \\
심사위원 가독성 저하 & 평가 누락 & Skills 로그 양방향 연결, 카드 하단 규칙 ID 배지 \\
데모 실패 & 심사 불가 & 샘플 CSV 파일 사전 제공, 원클릭 데모 옵션 \\
업로드 데이터 개인정보 & 공정성·저작권 이슈 & 공모전 데모에서는 공개 샘플 데이터 또는 합성 데이터를 사용한다. 실제 서비스 확장 시 업로드 파일은 세션 단위 임시 처리 후 삭제하는 것을 기본 정책으로 두며, 저장이 필요한 경우 사용자 동의와 보존 기간 고지를 별도로 요구한다 \\
\end{longtable}

\begin{center}\rule{0.5\linewidth}{0.5pt}\end{center}

\section{구현 로드맵}\label{uxad6cuxd604-uxb85cuxb4dcuxb9f5}

\subsection{Phase 1: 입력 처리 및 매핑 (핵심 진입 경로)}\label{phase-1-uxc785uxb825-uxcc98uxb9ac-uxbc0f-uxb9e4uxd551-uxd575uxc2ec-uxc9c4uxc785-uxacbduxb85c}

\begin{itemize}
\tightlist
\item
  파일 업로드 및 형식 검증 (F-UP-01)
\item
  데이터 미리보기 표시
\item
  AI 컬럼 매핑 + 신뢰도 기반 UI 분기 (F-CM-01, F-CM-02, F-CM-03)
\item
  분석 목적 입력 + 모호성 명확화 (F-GC-01, F-GC-02)
\item
  데이터 충분성 검사 (F-DS-01, F-DS-02)
\end{itemize}

\subsection{Phase 2: 분석 실행 및 대시보드 생성 (Skills.md 증명)}\label{phase-2-uxbd84uxc11d-uxc2e4uxd589-uxbc0f-uxb300uxc2dcuxbcf4uxb4dc-uxc0dduxc131-skills.md-uxc99duxba85}

\begin{itemize}
\tightlist
\item
  지표 계산 엔진 (F-AN-01\textasciitilde05)
\item
  대시보드 자동 구성 (F-DB-01)
\item
  AI 애널리스트 패널 + 금지 표현 검증 (F-AI-01)
\item
  Skills 로그 기록 및 열람 (F-SL-01)
\end{itemize}

\subsection{Phase 3: 데모 완성도 및 P1 경계 안내}\label{phase-3-uxb370uxbaa8-uxc644uxc131uxb3c4-uxbc0f-p1-uxacbduxacc4-uxc548uxb0b4}

\begin{itemize}
\tightlist
\item
  P0 데모 경로 최적화 (P0 샘플 데이터 기준, 30초 이내 완주 목표)
\item
  SCR-08 P1/P2 확장 안내 화면
\item
  시간 여유 시 Markdown 내보내기 (P1) 또는 벤치마크 비교(P1) 추가
\end{itemize}

\subsection{향후 확장 방향}\label{uxd5a5uxd6c4-uxd655uxc7a5-uxbc29uxd5a5}

본 MVP는 단일 CSV/XLSX 기반 포트폴리오 성과·리스크 대시보드 생성에 집중한다. 공모전 제출 이후에는 다음 방향으로 확장할 수 있다.

\begin{longtable}[]{@{}
  >{\raggedright\arraybackslash}p{(\columnwidth - 4\tabcolsep) * \real{0.3750}}
  >{\raggedright\arraybackslash}p{(\columnwidth - 4\tabcolsep) * \real{0.2500}}
  >{\raggedright\arraybackslash}p{(\columnwidth - 4\tabcolsep) * \real{0.3750}}@{}}
\toprule\noalign{}
\begin{minipage}[b]{\linewidth}\raggedright
확장 방향
\end{minipage} & \begin{minipage}[b]{\linewidth}\raggedright
설명
\end{minipage} & \begin{minipage}[b]{\linewidth}\raggedright
기대 효과
\end{minipage} \\
\midrule\noalign{}
\endhead
\bottomrule\noalign{}
\endlastfoot
Skills Pack & 성과 분석, 리스크 분석, 비중 분석, 벤치마크 분석 등 목적별 Skills.md 규칙 묶음 제공 & 분석 규칙 재사용성 강화 \\
Data Contract & 컬럼 매핑 결과와 데이터 충분성 판정 결과를 파일 구조별 계약으로 저장 & 반복 업로드 시 매핑 재사용, 분석 안정성 향상 \\
데이터 보완 추천 & 목적 달성에 필요한 추가 컬럼 또는 보조 파일 제안 & 데이터 부족 상황에서 사용자의 다음 행동 명확화 \\
Report Template Pack & 월간 성과 보고서, 리스크 점검 리포트 등 목적별 리포트 템플릿 제공 & 대시보드 결과의 보고서 전환 효율 향상 \\
외부 데이터 커넥터 & 벤치마크 수익률, 무위험수익률, 섹터 매핑, 종목 메타데이터 연동 & P1/P2 분석 범위 확대 \\
전략 명세 기반 대시보드 & 자연어 전략 설명을 구조화된 전략 명세로 변환하고 결과 대시보드 생성 & 퀀트 전략 분석 서비스로 확장 \\
Rule Library & 사용자 또는 팀 단위로 Skills.md 규칙을 저장·재사용 & 조직 내 분석 표준화 \\
\end{longtable}

\begin{center}\rule{0.5\linewidth}{0.5pt}\end{center}

\section{결론}\label{uxacb0uxb860}

LENS의 서비스 범위는 포트폴리오 CSV/XLSX 분석에 둔다. 범용 BI, 자유형 AI 챗봇, 투자 추천, 자동매매 기능은 본 MVP 범위에서 제외한다.

사용자가 투자 CSV/XLSX를 업로드하면 Skills.md에 정의된 금융 포트폴리오 분석 규칙이 컬럼 매핑 추론, 데이터 충분성 판정, 지표 계산, 시각화 선택, 인사이트 해석 템플릿, 금지 표현 검증을 순차적으로 처리한다. AI는 컬럼 추론 초기 제안과 자연어 해석 보조에 한정하며, 규칙 범위를 벗어난 판단을 수행하지 않는다.

본 기획서의 차별점은 자동화 수준보다 감사 가능성에 있다. Skills 로그에는 적용된 규칙 ID, 입력값, 계산식, 출력값, 분석 제외 사유가 기록된다. 심사위원은 대시보드 카드의 규칙 ID를 통해 분석 산출 과정을 직접 확인할 수 있다.

본 기획서는 Skills.md 규칙 명세, 처리 흐름 설계, Skills 로그 스키마의 세 층위에서 이 구조를 정의한다.


\end{document}
